% !TeX root = ../main.tex
% 12. fejezet — Összefoglalás és továbbfejlesztési lehetőségek
% A ftsrg/thesis-template-latex sablonba \include-dal illesztendő tartalomfájl.
% A fejezetre a 7., 8., 9. és 10. fejezet hivatkozik (chap:osszefoglalas) — a címke KÖTELEZŐ.
% Hivatkozott források: a dolgozat 5–11. fejezetei és a bennük dokumentált korlátok.

\chapter{Összefoglalás és továbbfejlesztési lehetőségek}
\label{chap:osszefoglalas}

\section{Eredmények}
\label{sec:ossz-eredmenyek}

A dolgozat egy teljes, működő webalkalmazás -- a HUNting vadászati közösségi és hirdetési
portál -- megtervezését és megvalósítását mutatta be, az adatbázistól a felhasználói felületig.
A bevezetésben kitűzött hat funkcióterület mindegyike megvalósult, a gerincfejezetekben
dokumentált terjedelemben. A \emph{felhasználókezelés} a kettős fióktípusra (egyéni vadász és
vadásztársaság) épülő regisztrációt, JWT-alapú hitelesítést, két szintű -- szerepkör- és
erőforrástulajdonos-alapú -- jogosultsági modellt és admin-felületet ad
(\ref{chap:felhasznalok}.~fejezet). A \emph{hirdetési piactér} kategorizált, szűrhető és
lapozott hirdetéslistát, képfeltöltéses hirdetéskezelést és kettős (felhasználói vagy
társasági) tulajdonlást valósít meg (\ref{chap:hirdetesek}.~fejezet). A \emph{társasági modul}
a tagfelvételi kérelmek teljes életciklusát, társaságon belüli szerepeket és tagoknak szóló
bejegyzéseket kezel (\ref{chap:tarsasagok}.~fejezet). A \emph{vadásznapló} képmelléklettel
rögzíti a személyes vadászati eseményeket (\ref{chap:naplo}.~fejezet). A \emph{térképes modul}
Google Maps-alapú megjelenítéssel, személyes és társasági rétegeken kezel jelölőket és
rajzolt területeket (\ref{chap:terkep}.~fejezet). Végül az \emph{AI-chatbot} szerveroldali
OpenAI-integrációval, szerep-szanitálással ad vadászati asszisztenst, az
\emph{időjárásoldal} pedig aktuális időjárási adatokat jelenít meg
(\ref{chap:chatbot}.~fejezet).

A megvalósítás mérnöki értékét a közös alapok adják: az egységes, 26 táblás, kifejezetten
megtervezett törlési viselkedésekkel és megszorításokkal bíró adatmodell, az idempotens
kezdeti adatfeltöltés, a következetes middleware-csővezeték és a kliensoldali közös
HTTP-réteg (\ref{chap:alapok}.~fejezet). A kódbázishoz háromszintű tesztkészlet készült --
124 backend-egységteszt, 17 komponenstesztfájl és 26 végponttól végpontig teszt --, amely a
jogosultsági ágakat, a külső integrációk izolálását és a kritikus felhasználói utakat fedi le
(\ref{chap:teszteles}.~fejezet). A dolgozat végig törekedett az őszinte minősítésre: minden
funkciófejezet külön szakaszban rögzítette a megvalósítás korlátait, megkülönböztetve a
működő, az előkészített és a hiányzó képességeket.

\section{Továbbfejlesztési lehetőségek}
\label{sec:ossz-tovabbfejlesztes}

A funkciófejezetekben dokumentált korlátok egyben a továbbfejlesztés térképét is kirajzolják.
A legfontosabb irányok a következők.

\begin{itemize}
	\item \emph{Hírfolyam befejezése:} a backend-oldalon előkészített hírmodul
	      (\ref{chap:tarsasagok}.~fejezet) felhasználói felületének elkészítése, a hiányzó
	      autorizáció pótlásával együtt -- ez a rendszer legégetőbb biztonsági adóssága.
	\item \emph{Hitelesítés megerősítése:} frissítő token és tokenvisszavonás bevezetése, a
	      jelszószabályzat szigorítása, a konfigurációs titkok (JWT-kulcs, API-kulcsok)
	      titokkezelőbe szervezése.
	\item \emph{Térképi funkciók kiteljesítése:} az adatmodell szintjén előkészített útvonalak
	      (\texttt{MapPath}) végpontjainak és felületének elkészítése; a láthatósági kódok
	      tényleges szűrési logikájának megvalósítása; szerkesztőfelület a meglévő
	      PUT-végpontokhoz (\ref{chap:terkep}.~fejezet).
	\item \emph{Felületi hiányok pótlása:} szerkesztő- és törlőműveletek a vadásznaplóban,
	      lapozás a naplólistában (\ref{chap:naplo}.~fejezet), admin-felület a létrehozó és
	      módosító végpontokhoz, valamint a felületi szövegek ékezeteinek rendezése.
	\item \emph{Külső integrációk keményítése:} az anonim chatvégpont sebességkorlátozása és a
	      kliensoldali OpenWeather-hívás backend-proxy mögé szervezése
	      (\ref{chap:chatbot}.~fejezet).
	\item \emph{Közösségi funkciók bővítése:} értesítések (például tagfelvételi döntésről),
	      tagok közötti üzenetküldés, valamint hosszabb távon natív mobilalkalmazás.
	\item \emph{Üzemeltetési érettség:} CI-futtatás a meglévő tesztkészletre, konténerizálás,
	      környezetenkénti konfiguráció és telepítési folyamat; kódszinten a lusta
	      oldalbetöltés bekapcsolása és a kliensoldali típusok generálása az
	      OpenAPI-leírásból.
\end{itemize}

A felsorolás terjedelme nem a rendszer hiányosságát, hanem a választott megközelítés
következményét tükrözi: a dolgozat egy szélesen integrált rendszert valósított meg hat
funkcióterületen, teljes vertikumban, és minden területen pontosan dokumentálta, hol húzódik a
megvalósítás határa. Az elkészült alap -- a közös adatmodell, a jogosultsági rendszer és a
tesztkészlet -- úgy lett kialakítva, hogy ezek az irányok töréspont nélkül építhetők rá.
